% ============================================================
% MSc Dissertation — Maria Bautista, University of Essex (CE902)
% Inductive Graph Neural Networks for Cold-Start Bus Demand
% ============================================================
\documentclass[12pt,a4paper]{report}

\usepackage[utf8]{inputenc}
\usepackage[T1]{fontenc}
\usepackage{lmodern}
\usepackage[margin=2.5cm]{geometry}
\usepackage{amsmath,amssymb}
\usepackage{graphicx}
\usepackage{booktabs}
\usepackage{caption}
\usepackage{subcaption}
\usepackage{hyperref}
\usepackage{cleveref}
\usepackage{setspace}
\usepackage{natbib}
\onehalfspacing

\hypersetup{colorlinks=true, linkcolor=black, citecolor=blue, urlcolor=blue}

\title{\textbf{Inductive Graph Neural Networks for Cold-Start Prediction of Bus-Stop Passenger Demand}}
\author{Maria Bautista}
\date{Supervisor: Dr Vishal K. Singh \\ University of Essex \\ \today}

\begin{document}

\maketitle

% ============================================================
\begin{abstract}
\noindent
Predicting passenger demand at bus stops that have no prior ridership history
is a practical challenge facing urban transport planners whenever new
infrastructure is developed. Conventional spatio-temporal models rely on
historical flow data and are unable to produce estimates for stops that were
absent during training---a limitation referred to here as the
\emph{urban cold-start problem}. This dissertation investigates whether an
inductive Graph Attention Network (GATv2) can predict bus-stop boardings for
geographically unseen areas using only land-use accessibility features and
network topology, without any historical demand at the target locations.

A multigraph is constructed over 17,943 London bus stops by combining
geographic proximity edges (K-nearest neighbours, $K{=}5$) with route-connectivity
edges derived from Transport for London's BUSTO survey. Node features are drawn
from the AI23 accessibility indicators, measuring 30-minute public transport
access to employment, healthcare, education, and retail at LSOA level.
Performance is evaluated under a 33-fold leave-borough-out cross-validation,
where all stops in one London borough are held out at each fold to simulate
the arrival of a geographically new area.

GATv2 achieves a Weighted Mean Absolute Percentage Error (WMAPE) of 0.838 across
33 folds, compared with 0.812 for a non-graph Random Forest and 1.082 for a
Historical Average baseline. GATv2 outperforms the Random Forest in 7 of 33
boroughs, consistently those with dense route networks in inner London, where
message passing over route-connectivity edges provides signal beyond what
tabular features alone carry. An MLP ablation (WMAPE 0.814) confirms that
node-level accessibility features account for most of the predictive power, with
graph structure contributing additional but bounded benefit. The residual gap
between GATv2 and the Random Forest is attributed to feature homophily arising
from LSOA-level aggregation, and stop-level POI features are identified as the
most impactful direction for future work.
\end{abstract}

\tableofcontents
\listoffigures
\listoftables

% ============================================================
\chapter{Introduction}
\label{ch:intro}

\section{Background and motivation}
Urban transit authorities now collect more data than ever before: smart-card
records, vehicle location traces, real-time sensor feeds, and detailed land-use
registries. Yet the methodological frameworks most authorities rely on were
designed for cities that change slowly. London is a stark example. The city
adds residential developments, repurposes industrial land, and extends rail
lines at a pace no static demand model was built to track. The opening of the
Elizabeth Line in 2022 and the residential expansion into areas such as
Barking Riverside illustrate the difficulty directly: planners must estimate
demand for places that have little or no prior travel history at precisely the
moment when accurate estimates would be most valuable.

This is the problem the present work addresses, and which we refer to as the
\emph{urban cold-start}. The term is borrowed from recommender systems, where
it denotes the difficulty of making predictions for a user or item with no
interaction history \citep{schein2002methods}. Transposed to urban geography,
it describes a transit node---a new stop, a newly served neighbourhood---for
which no historical passenger record exists. Conventional forecasting tools
derive their predictive power from historical flow patterns and are therefore,
by construction, blind to such locations.

Bus line networks are complex, nonlinear, time-varying systems, and modelling
the spatio-temporal correlations of passenger flows across a city is a
long-standing challenge \citep{zheng2025mfstgat}. Most recent advances have
improved accuracy substantially on \emph{known} networks, yet they share a
common assumption: that the set of stops observed during training is the same
set that will be queried at inference time. For a transport authority managing
a growing city, this assumption fails repeatedly.

\section{The research gap}
A critical limitation runs through the recent literature on spatio-temporal
graph neural networks for transit demand: these models are predominantly
\emph{transductive}. A transductive model is trained on a fixed graph of nodes
and edges and has no principled mechanism for handling nodes absent during
training \citep{hamilton2017inductive}. Such models are highly effective at
describing demand within an existing, immutable network, but they cannot reason
about infrastructure that does not yet exist without full retraining from
scratch. They describe what is; they cannot anticipate what is not yet there.

The alternative pursued here is \emph{inductive} learning. Rather than learning
node-specific parameters tied to a particular graph, an inductive model learns
a general mapping from a node's features and local neighbourhood structure to a
prediction \citep{hamilton2017inductive}. Such a model can, in principle,
produce an estimate for a node it has never seen, provided that node's features
and immediate surroundings are available. This reframing---from memorising a
fixed network to learning a transferable function over local graph
structure---is what makes cold-start prediction tractable, and it is the
central idea of this dissertation.

\section{Aim and research questions}
The aim of this dissertation is to determine whether an inductive graph neural
network can infer bus-stop passenger demand for areas with no prior travel
history, using only the land-use accessibility profile and route-connectivity
structure of their surroundings. To evaluate this, a \emph{leave-borough-out}
protocol is adopted: entire London boroughs are withheld from training and the
model is required to predict demand at their stops as though those areas had
never been observed. This simulates the arrival of a new, unseen district.
The investigation is structured around one primary research question and one
ablation sub-question:

\begin{itemize}
  \item \textbf{RQ1.} Under a leave-borough-out evaluation that simulates
  geographically unseen areas, can an inductive graph attention network (GATv2)
  predict bus-stop passenger demand more accurately than (a) a
  Historical-Average baseline and (b) a non-graph Random Forest trained on the
  same features?
  \item \textbf{RQ1a (ablation).} How much of any predictive improvement is
  attributable to the \emph{graph structure} itself, as opposed to the node
  features alone?
\end{itemize}

The ablation sub-question is the methodological core of the work. Demonstrating
that a graph model outperforms a baseline is uninformative unless one can
establish that the relational structure, rather than merely a richer feature
set, is responsible for the gain.

\section{Contributions}
This dissertation makes three contributions. First, it reformulates bus stop
demand prediction as an inductive cold-start problem, enabling estimates for
stops in areas never seen during training. Second, it adopts a spatially honest
evaluation---leave-borough-out cross-validation---that isolates the contribution
of graph structure from that of node features through an MLP ablation, addressing
a question that predictive comparisons frequently leave unanswered. Third, it
delivers a fully reproducible pipeline built on openly available United Kingdom
data sources---TfL BUSTO boardings, AI23 accessibility indicators, and ONS
boundary data---allowing the analysis to be replicated and extended without
proprietary access.

\section{Dissertation structure}
Chapter~\ref{ch:litreview} reviews spatio-temporal graph models, inductive
representation learning, and land-use as a predictor of mobility, positioning
the present work within this literature. Chapter~\ref{ch:data} describes the
datasets and the construction of the multigraph. Chapter~\ref{ch:method} sets
out the inductive GATv2 model, the baselines, the validation design, and the
evaluation metrics. Chapter~\ref{ch:results} reports the predictive results and
the structural ablation. Chapter~\ref{ch:discussion} interprets the findings,
and Chapter~\ref{ch:conclusion} concludes with limitations and directions for
future work.


% ============================================================
\chapter{Literature Review}
\label{ch:litreview}

\section{Spatio-temporal graph models for transit demand}
Graph Neural Networks (GNNs) have become the dominant paradigm for transit
demand modelling over the past five years, largely because passenger flows
are inherently relational: demand at one stop is not independent of demand
at nearby stops or stops on shared routes. The foundational work of
\citet{kipf2017semi} introduced the Graph Convolutional Network (GCN), which
performs localised spectral convolution over a graph by aggregating features
from each node's immediate neighbours with equal weight. Applied to traffic
and transit networks, GCN-based models learn representations that reflect
both a node's own features and the characteristics of its surroundings, and
they have been shown to outperform independent-node models on speed and flow
prediction tasks \citep{kipf2017semi}.

A key limitation of GCN for transit demand is the assumption that all
neighbours contribute equally to a node's representation. In a bus network,
a stop at a busy interchange should receive much stronger influence from its
high-demand neighbours than from a low-frequency stop that happens to be
geographically nearby. \citet{velickovic2018graph} addressed this by
introducing the Graph Attention Network (GAT), which assigns learnable
attention weights to each neighbour and aggregates accordingly. However,
\citet{brody2022how} subsequently demonstrated that the attention mechanism
in GAT is mathematically \emph{static}: the attention score between nodes $i$
and $j$ can be decomposed into additive terms that depend only on $i$ or
only on $j$, meaning the model cannot express attention that is jointly
conditioned on the pair. They proposed GATv2, which introduces a single
architectural change---applying the weight matrix to the concatenation of
both node features before the non-linearity---that makes attention genuinely
dynamic. This distinction is particularly relevant for heterogeneous urban
graphs where edge importance varies with context, and GATv2 is the model
adopted in this work.

For bus demand specifically, \citet{zheng2025mfstgat} proposed the
Multigraph Fusion Spatio-Temporal Graph Attention Network (MF-STGAT), which
fuses spatial proximity graphs with semantic functional graphs and incorporates
temporal attention to model the time-varying nature of passenger flows at
the regional level. MF-STGAT achieves strong predictive accuracy on observed
stops and represents the current state of the art for this task. However, it
is a transductive model: the learned node embeddings are tied to specific
stops and cannot be transferred to stops absent from the training graph.
This is the gap that the present dissertation addresses.

\section{Inductive representation learning on graphs}
The distinction between transductive and inductive learning is central to
the cold-start problem. Transductive models, such as standard GCN and most
spatio-temporal GNNs for transit, optimise embeddings for a fixed set of
nodes. When the graph changes---when new nodes arrive---these embeddings must
be recomputed or the model retrained. This is computationally expensive and
operationally impractical for a transit authority that adds stops continuously.

\citet{hamilton2017inductive} introduced GraphSAGE, the first inductive
GNN framework, and drew the distinction explicitly. GraphSAGE learns
\emph{aggregation functions}---rather than fixed embeddings---that can be
applied to any node given its features and neighbourhood. At inference time,
a new node's embedding is computed by passing its features and those of its
neighbours through the learned aggregators, without retraining. This is
precisely the operation required for cold-start prediction: the model never
sees the test stops' demand labels but can still produce predictions from
their features and network context.

GATv2 \citep{brody2022how} is inductive in the same sense. The attention
parameters are shared across all node pairs, not tied to specific node
identities. A new node's representation is computed by applying the same
attention function to its features and those of its neighbours. In the present
work, this is implemented using PyTorch Geometric's \texttt{subgraph} utility
to extract the training subgraph at each fold, ensuring that test borough
stops are processed inductively during evaluation.

\section{Land-use and accessibility as predictors of mobility}
A complementary literature has examined how well land-use and accessibility
features can predict travel demand in the absence of historical flow data.
\citet{simini2021deep} demonstrated that a deep learning model fed with
population and point-of-interest data could reproduce observed mobility flows
across multiple cities without any prior trip data---a result with direct
implications for the cold-start setting. The intuition is that the built
environment shapes the latent demand for travel: an area with dense employment
and poor residential character will generate outbound commute demand, while a
predominantly residential area will generate inbound demand in the same period.

The AI23 Accessibility Indicators \citep{verduzco2024ai23} operationalise
this idea at the LSOA level for Great Britain, computing the number of key
destinations---employment sites, hospitals, GP practices, supermarkets,
pharmacies, and schools---reachable within 30 minutes by public transport
during the AM peak. These indicators have been validated as correlates of
transport usage in the literature and represent the primary feature set in
this work. Their LSOA-level granularity, while a limitation discussed further
in Section~\ref{sec:limitations}, reflects the scale at which socioeconomic
and land-use data are most reliably available in the United Kingdom.

\section{The cold-start problem and spatial validation}
The cold-start problem as studied here is an instance of a more general
challenge in spatial machine learning: the tendency of standard evaluation
protocols to overestimate the generalisation ability of models trained on
spatially autocorrelated data. \citet{roberts2017cross} showed that randomly
partitioning spatially structured data into training and test folds allows
nearby observations to appear in both splits, violating the independence
assumption underlying cross-validation and producing unrealistically
optimistic error estimates. Their recommendation---to use spatially blocked
folds, such that training and test sets are separated by a minimum geographic
distance or administrative boundary---has become standard practice in
ecological and geographic modelling. The leave-borough-out design adopted
here follows this recommendation directly: the administrative boundary of
each London borough serves as the blocking criterion, guaranteeing that no
training stop is geographically adjacent to a test stop within the same fold.

\section{Research gap and positioning}
Table~\ref{tab:litgap} summarises the positioning of this work relative to the
key prior studies. The critical distinguishing characteristic is the
combination of inductive learning, spatial cold-start evaluation, and the
explicit ablation of graph structure versus node features. None of the
reviewed studies simultaneously satisfy all three criteria. MF-STGAT achieves
strong performance but is transductive. GraphSAGE is inductive but was not
applied to transit demand or spatially blocked evaluation. The present work
occupies the intersection: an inductive GNN evaluated under conditions that
genuinely simulate the deployment scenario it targets.

\begin{table}[ht]
  \centering
  \caption{Positioning of this work relative to key prior studies.}
  \label{tab:litgap}
  \begin{tabular}{lcccc}
    \toprule
    Study & Inductive & Transit demand & Spatial CV & Graph ablation \\
    \midrule
    \citet{kipf2017semi}          & No  & No  & No  & No  \\
    \citet{velickovic2018graph}   & No  & No  & No  & No  \\
    \citet{hamilton2017inductive} & Yes & No  & No  & No  \\
    \citet{brody2022how}          & Yes & No  & No  & No  \\
    \citet{zheng2025mfstgat}      & No  & Yes & No  & No  \\
    \textbf{This work}            & \textbf{Yes} & \textbf{Yes} & \textbf{Yes} & \textbf{Yes} \\
    \bottomrule
  \end{tabular}
\end{table}


% ============================================================
\chapter{Data and Multigraph Construction}
\label{ch:data}

\section{Overview}
Three datasets are used in this study: (1) TfL BUSTO, which provides the
prediction target; (2) TfL Bus Stop Locations, which provides geographic
coordinates; and (3) the AI23 Accessibility Indicators, which provide node
features. All three datasets are openly available under the TfL Open Data
Licence or the Open Government Licence and can be obtained without special
access arrangements. The processing pipeline, implemented in four sequential
Python scripts, transforms these raw sources into the 17,943-node multigraph
used for modelling.

\section{Target variable: TfL BUSTO boardings}
The Bus Stop Usage and Occupancy Survey (BUSTO; TfL, 2023/24) provides
estimated weekday passenger boardings for every bus stop served by the TfL
network, broken down by route, direction, and 15-minute time interval. Each
row in the raw data corresponds to one stop on one route in one direction
during one quarter-hour period, with fields: \texttt{YEAR}, \texttt{DAY\_TYPE},
\texttt{TIMEBAND}, \texttt{QHr}, \texttt{ROUTE}, \texttt{DIRECTION},
\texttt{STOPCODE}, \texttt{STOPNAME}, \texttt{STOPSEQUENCE},
\texttt{Boardings}, \texttt{Alightings}, \texttt{Load}, \texttt{Capacity},
\texttt{Seats}, \texttt{V/C}. The four BUSTO CSV files together contain
3,441,745 rows.

The target variable \texttt{total\_boardings} is constructed by summing
\texttt{Boardings} across all weekday rows for each unique \texttt{STOPCODE}.
Weekend and public holiday records are excluded, as weekday demand is
substantially higher and more representative of the commute-driven travel
patterns that accessibility features are designed to capture. The resulting
stop-level dataset contains 19,682 unique stops before geographic filtering.
The target is right-skewed: the median stop records approximately 160 boardings
per typical weekday, while the busiest stops exceed 5,000. A $\log(1 + y)$
transformation is applied to the target during model training and inverted for
metric computation.

It is important to note that BUSTO boardings are survey-derived estimates
rather than exact counts. Vehicle occupancy and alightings are algorithmically
inferred from ticketing tap data and on-board surveys; the fractional values
observed (e.g., 968.09 for Westminster Stn / Parliament Square) reflect
statistical weighting applied during this estimation process. Boardings are
retained as the target because they are the most directly observed quantity in
the survey.

\section{Geographic coordinates: TfL Bus Stop Locations}
The TfL Bus Stop Locations file (\texttt{Bus\_Stops.csv}) provides, for each
stop, a unique \texttt{STOP\_CODE} identifier and coordinates in the British
National Grid (EPSG:27700) projection. A matching verification confirmed 99.6\%
overlap between BUSTO \texttt{STOPCODE} values and the locations file
\texttt{STOP\_CODE} column. Easting and Northing coordinates were reprojected
to WGS84 latitude/longitude (EPSG:4326) using pyproj \citep{snow2023pyproj}.
Stops for which no coordinates could be matched (0.4\%) were excluded as
non-physical or virtual stops. After this join, 19,612 stops retained
coordinates.

\section{Node features: AI23 Accessibility Indicators}
\label{sec:ai23}
The AI23 Accessibility Indicators \citep{verduzco2024ai23} provide, for each
2011 LSOA in England, a set of accessibility scores measuring the number of
key destinations reachable within specified travel time thresholds by public
transport, under different temporal conditions. This study extracts eight
features, all measured at the 30-minute threshold under public transport
(\texttt{mode = pt}) during the AM peak (\texttt{time\_of\_day = am}):

\begin{enumerate}
  \item \texttt{employment\_all\_30min} --- total employment sites accessible
  \item \texttt{hospitals\_30min} --- hospitals and major health centres
  \item \texttt{gp\_30min} --- GP practices
  \item \texttt{supermarkets\_30min} --- major supermarkets
  \item \texttt{pharmacies\_30min} --- pharmacies
  \item \texttt{primary\_schools\_30min} --- primary schools
  \item \texttt{secondary\_schools\_30min} --- secondary schools
  \item \texttt{main\_bua\_30min} --- major built-up areas (urban centres)
\end{enumerate}

The AM peak filter was chosen to align with the temporal profile of BUSTO
weekday demand, which is dominated by morning commute trips. All eight features
exhibit right-skewed distributions (skewness $> 1.7$) and are $\log(1+x)$
transformed before model input.

Each stop was assigned a 2011 LSOA code via two steps. First, its postcode
(obtained from the TfL Bus Stops file) was submitted to the postcodes.io batch
API \citep{postcodesio}, which returns LSOA code and London borough for each
postcode. Second, because postcodes.io returns 2021 LSOA codes for areas
subject to the 2021 census boundary revision, and AI23 uses 2011 geography,
the ONS LSOA 2021-to-2011 best-fit crosswalk was applied to recover unmatched
stops \citep{ons2022crosswalk}. This recovered 555 of 1,123 initially
unmatched stops. The remaining 568 stops belong to genuine 2021-only LSOAs
with no single 2011 parent and were excluded (3.1\% of London stops).

After filtering to the 33 London boroughs (LAD codes beginning \texttt{E09}),
the final dataset contains \textbf{17,943 stops} across 33 boroughs. A notable
characteristic of the resulting feature matrix is that 96.4\% of stop pairs
within the same LSOA share identical AI23 feature vectors, because features
are assigned at LSOA level rather than individual stop level. This has
important implications for message passing, discussed in
Section~\ref{sec:homophily}.

\section{Multigraph construction}
\label{sec:multigraph}
The transit network is represented as a directed multigraph $G = (V, E)$ where
$V$ is the set of 17,943 bus stops and $E$ is a union of two edge sets with
different semantic meanings.

\subsection{Geographic KNN edges}
For each stop $v_i$, edges are created to its $K=5$ nearest stops by Haversine
great-circle distance. Haversine distance is used rather than Euclidean
distance to account for the curvature of the Earth over London's spatial
extent. A BallTree data structure with the Haversine metric was used for
efficient computation \citep{scikit2011}. All edges are made bidirectional.
This produces \textbf{109,182 directed edges}. The role of geographic edges is
to propagate spatial context: stops that are physically close tend to serve
areas with similar demographic and accessibility profiles, and their demand
should be correlated.

\subsection{Route-connectivity edges}
For each bus route and direction in the BUSTO dataset, stops are ordered by
their \texttt{STOPSEQUENCE} value and directed edges are created between
consecutive stops in both directions. These edges capture functional
network flow: passengers who board at stop $v_i$ on route $r$ in direction $d$
generate correlated demand at the next stop $v_{i+1}$ on the same service,
because they are part of the same through-journey. This produces
\textbf{43,220 directed edges} after deduplication.

The rationale for combining both edge types is that neither is sufficient alone.
Geographic edges connect stops that may be physically close but functionally
unrelated (e.g., stops on parallel routes with different catchment areas).
Route edges without geographic context miss neighbourhood-level demand
regularities that influence all stops in an area regardless of route membership.
After deduplication, the combined multigraph contains
\textbf{124,234 directed edges}.

\subsection{Node feature augmentation}
Because 96.4\% of stops within the same LSOA share identical AI23 feature
vectors, geographic coordinates (latitude and longitude) are appended to the
feature matrix, giving each stop a unique geometric position even when its
accessibility indicators are shared with nearby stops. The final input feature
dimension is therefore 10: eight log-transformed AI23 indicators plus latitude
and longitude. This design decision follows \citet{zheng2025mfstgat}, who use
stop-level POI features to ensure node uniqueness; geographic coordinates serve
as an approximation where stop-level POI data is unavailable.


% ============================================================
\chapter{Methodology}
\label{ch:method}

\section{Problem formulation}
Let $G = (V, E, X)$ denote the multigraph defined in Chapter~\ref{ch:data},
where $X \in \mathbb{R}^{|V| \times 10}$ is the node feature matrix and
$y \in \mathbb{R}^{|V|}$ is the vector of total weekday boardings. The
cold-start prediction task is defined as follows: given a partition
$V = V_{\text{train}} \cup V_{\text{test}}$ with $V_{\text{train}} \cap V_{\text{test}} = \emptyset$,
learn a function $f: (X_{V_{\text{train}}}, E_{\text{train}}) \rightarrow \hat{y}_{V_{\text{test}}}$
such that the model can produce accurate predictions for test nodes
\emph{without access to} $y_{V_{\text{test}}}$ during training.
Crucially, the test nodes correspond to an entire geographic region
(a London borough) that was completely absent from training, not a random
sample of individual stops.

\section{Inductive GATv2}
\label{sec:gatv2}
Graph Attention Network version 2 (GATv2; \citealt{brody2022how}) is chosen
as the primary model for three reasons. First, it is inductive: its parameters
are functions over node features and neighbourhood structure, not tied to
specific node identities, so it can produce embeddings for nodes never seen
during training. Second, its dynamic attention mechanism is strictly more
expressive than both GCN (which uses fixed normalised weights) and the
original GAT (whose attention is mathematically static). Third, it achieves
strong performance on graph regression tasks with heterogeneous node features.

The unnormalised attention score between nodes $i$ and $j$ in GATv2 is:
\begin{equation}
  e_{ij} = \vec{a}^{\top}\,\mathrm{LeakyReLU}\!\left(W\,[\vec{h}_i \,\Vert\, \vec{h}_j]\right),
  \label{eq:gatv2attn}
\end{equation}
where $W \in \mathbb{R}^{d' \times 2d}$ is a shared weight matrix,
$\vec{a} \in \mathbb{R}^{d'}$ is a learnable attention vector,
$\vec{h}_i$ and $\vec{h}_j$ are the current node representations, and
$\Vert$ denotes concatenation. Scores are normalised across each node's
neighbourhood using softmax:
\begin{equation}
  \alpha_{ij} = \frac{\exp(e_{ij})}{\sum_{k \in \mathcal{N}(i)} \exp(e_{ik})},
\end{equation}
and the updated representation of node $i$ is:
\begin{equation}
  \vec{h}'_i = \sigma\!\left(\sum_{j \in \mathcal{N}(i)} \alpha_{ij}\,W'\,\vec{h}_j\right).
\end{equation}

The architecture used in this study consists of two GATv2 convolutional layers.
The first layer uses $H=4$ attention heads with hidden dimension $d=64$ per
head (output dimension $4 \times 64 = 256$); the second layer uses a single
head with dimension 64, followed by a linear readout layer producing a scalar
prediction. ELU activation and dropout ($p=0.3$) are applied after the first
layer. The full feature dimension at input is 10, matching the 8 log-transformed
AI23 indicators plus latitude and longitude.

\subsection{True inductive evaluation}
Following \citet{hamilton2017inductive}, the inductive setting is enforced by
extracting a subgraph of training nodes at each fold using PyTorch Geometric's
\texttt{subgraph} utility. Test borough stops are added to this subgraph only
at evaluation time, and their neighbourhood messages are aggregated from
training nodes through the learned attention function. The model never receives
the demand labels of test nodes during training in any fold.

\section{Training protocol}
The model is trained by minimising mean squared error on $\log(1+y)$:
\begin{equation}
  \mathcal{L} = \frac{1}{|V_{\text{train}}|} \sum_{i \in V_{\text{train}}}
  \left(\log(1 + \hat{y}_i) - \log(1 + y_i)\right)^2.
\end{equation}

Optimisation uses Adam \citep{kingma2015adam} with learning rate $10^{-3}$.
Training runs for a maximum of 300 epochs with early stopping: training halts
when validation loss (computed on a random 10\% of training stops) does not
improve by more than $10^{-5}$ for 40 consecutive epochs. The best checkpoint
by validation loss is restored before evaluation.

Predictions are converted back to the original scale via $\hat{y} = e^{\hat{z}} - 1$
and clamped to the interval $[0,\, \log(1 + 2 \cdot \max(y_{\text{train}}))]$
before inversion, preventing extreme extrapolation on boroughs with
out-of-distribution demand profiles.

\section{Baselines}
Three baselines are evaluated alongside GATv2.

\paragraph{Historical Average (HistAvg).}
The simplest conceivable baseline: predict every test stop's demand as the
mean boardings of all training stops. This baseline uses no features and no
graph structure. It establishes the minimum bar; any informative model should
surpass it.

\paragraph{Random Forest (RF).}
An ensemble of 300 decision trees trained on the 10-dimensional standardised
feature matrix, predicting $\log(1+y)$ and converting back at evaluation.
RF is the primary non-graph comparator and, crucially, uses exactly the same
features as GATv2. Any difference between RF and GATv2 performance can
therefore be attributed to the graph structure, not to feature differences.
This was trained with \texttt{n\_jobs=-1} and \texttt{random\_state=42}.

\paragraph{Multi-Layer Perceptron (MLP, ablation).}
A two-hidden-layer neural network (64 units, ReLU, dropout 0.3) trained on
the same 10 features with the same Adam optimiser, learning rate, and early
stopping as GATv2, but without any graph structure. The MLP ablation answers
RQ1a directly: by comparing MLP to GATv2, one isolates the contribution of
graph structure while controlling for the neural architecture.

\section{Spatial cross-validation}
\label{sec:spatialcv}
London's 33 administrative boroughs serve as the blocking units for
leave-borough-out cross-validation. At fold $k$, all stops in borough $b_k$
are held out as the test set; the remaining 32 boroughs form the training set.
This is repeated for all 33 boroughs, yielding 33 folds with test sizes
ranging from 101 stops (City of London) to 1,179 stops (Bromley).

A StandardScaler is fitted exclusively on training stops at each fold and
applied to both training and test features. Fitting the scaler on the full
dataset before splitting would constitute data leakage, as the scaling
statistics of the test borough would influence the model's input representation
\citep{kaufman2012leakage}.

\section{Evaluation metrics}
Three metrics are reported. WMAPE is the primary metric because it is
robust to stops with near-zero demand, where plain MAPE is undefined or
arbitrarily large:
\begin{equation}
  \mathrm{WMAPE} = \frac{\sum_i |\hat{y}_i - y_i|}{\sum_i y_i + \epsilon}.
  \label{eq:wmape}
\end{equation}
RMSE and MAE are reported for completeness. All metrics are computed on the
original (uninverted) demand scale. Results are summarised as mean $\pm$
standard deviation and median across 33 folds.

\section{Implementation}
The full pipeline is implemented in Python 3.11 using PyTorch 2.12 and PyTorch
Geometric 2.8.0 \citep{fey2019fast}. Scikit-learn 1.5 \citep{scikit2011} is
used for Random Forest and StandardScaler. All experiments were run on a
CPU-only environment (Intel Core i7, 16 GB RAM); total wall-clock time for the
33-fold run was 35 minutes. A fixed random seed (\texttt{seed=42}) was used
for all stochastic components. The complete code and processed dataset are
provided in the supplementary material.


% ============================================================
\chapter{Experiments and Results}
\label{ch:results}

\section{Experimental setup}
The multigraph described in Chapter~\ref{ch:data} contains 17,943 nodes,
124,234 directed edges, and a 10-dimensional node feature matrix. For each of
the 33 leave-borough-out folds, the model is trained from scratch on the
subgraph induced by the 32 training boroughs and evaluated on the held-out
borough. Hyperparameters were set prior to the full evaluation run based on
alignment with \citet{zheng2025mfstgat} (hidden dimension 64, 4 attention
heads) and are not tuned per fold. A five-fold quick run on the first five
boroughs alphabetically was used to verify the pipeline before committing to
the full 33-fold experiment.

\section{Overall predictive performance (RQ1)}
Table~\ref{tab:results} reports WMAPE, RMSE, and MAE for all four models,
averaged across the 33 folds.

\begin{table}[ht]
  \centering
  \caption{Cold-start prediction performance: 33-fold leave-borough-out CV.
           Mean $\pm$ std (median) across boroughs.}
  \label{tab:results}
  \begin{tabular}{lcccccc}
    \toprule
    Model & WMAPE mean & WMAPE std & WMAPE med & RMSE mean & MAE mean \\
    \midrule
    Historical Average       & 1.0822 & 0.3778 & 0.9636 & 554.0 & 316.1 \\
    Random Forest (no graph) & 0.8121 & 0.0258 & 0.8078 & 561.4 & 264.4 \\
    MLP (ablation)           & 0.8136 & 0.0309 & 0.8082 & 561.6 & 264.2 \\
    GATv2 (inductive)        & 0.8381 & 0.0321 & 0.8378 & 576.8 & 272.6 \\
    \bottomrule
  \end{tabular}
\end{table}

All three learned models substantially outperform the Historical Average
baseline (WMAPE 1.082), confirming that both accessibility features and graph
structure carry meaningful signal beyond a simple mean prediction. The
Random Forest achieves WMAPE 0.812, GATv2 achieves 0.838, and the MLP
ablation achieves 0.814. GATv2 reduces WMAPE by 0.244 relative to HistAvg
but trails RF by 0.026. The narrow gap between RF and MLP (0.002) indicates
that the two models extract broadly equivalent information from the same
feature set, which is expected given their similar functional capacity on
low-dimensional inputs.

\section{The value of graph structure (RQ1a)}
The MLP ablation (WMAPE 0.814) is nearly identical to Random Forest (0.812)
and both substantially below GATv2 (0.838). This finding is central to the
interpretation of the dissertation: the accessibility features alone account
for almost all recoverable predictive signal in this dataset, and the graph
structure adds only a bounded increment.

The explanation lies in the feature homophily identified in
Section~\ref{sec:ai23}: 96.4\% of stops share their LSOA with at least one
other stop, meaning they have identical AI23 feature vectors. When GATv2 applies
message passing to aggregate from neighbours, it is averaging representations
that are nearly identical to the node's own, so the aggregated message contains
little new information. The addition of latitude and longitude partially
mitigates this---each stop now has a unique geometric position even when its
accessibility indicators match its neighbours---but does not fully substitute
for stop-level unique features.
\label{sec:homophily}

\section{Borough-level analysis}
\label{sec:borough}
While aggregate performance favours RF, GATv2 outperforms RF in 7 of 33
boroughs, and ties in one (Enfield, WMAPE 0.860 for both). Table~\ref{tab:borough}
presents the full per-borough results.

\begin{table}[ht]
  \centering
  \caption{Per-borough WMAPE results for all four models.
           GATv2 wins are highlighted in bold.}
  \label{tab:borough}
  \small
  \begin{tabular}{lrcccc}
    \toprule
    Borough & $n$ & HistAvg & RF & MLP & GATv2 \\
    \midrule
    Barking and Dagenham   & 364  & 0.893 & 0.778 & 0.777 & 0.878 \\
    \textbf{Barnet}        & 993  & 1.240 & 0.853 & 0.837 & \textbf{0.830} \\
    Bexley                 & 672  & 1.697 & 0.825 & 0.828 & 0.868 \\
    Brent                  & 575  & 0.818 & 0.797 & 0.798 & 0.821 \\
    \textbf{Bromley}       & 1179 & 2.607 & 0.872 & 0.882 & \textbf{0.870} \\
    Camden                 & 424  & 0.867 & 0.802 & 0.787 & 0.820 \\
    \textbf{City of London}& 101  & 0.798 & 0.819 & 0.798 & \textbf{0.798} \\
    Croydon                & 1019 & 1.548 & 0.789 & 0.793 & 0.815 \\
    Ealing                 & 683  & 0.964 & 0.827 & 0.830 & 0.868 \\
    Enfield                & 707  & 1.145 & 0.860 & 0.855 & 0.860 \\
    Greenwich              & 659  & 0.999 & 0.831 & 0.833 & 0.894 \\
    \textbf{Hackney}       & 442  & 0.827 & 0.789 & 0.777 & \textbf{0.767} \\
    Hammersmith and Fulham & 276  & 0.788 & 0.799 & 0.794 & 0.862 \\
    \textbf{Haringey}      & 461  & 0.840 & 0.806 & 0.803 & \textbf{0.801} \\
    Harrow                 & 432  & 1.088 & 0.794 & 0.792 & 0.803 \\
    Havering               & 691  & 1.602 & 0.846 & 0.854 & 0.866 \\
    Hillingdon             & 706  & 1.300 & 0.856 & 0.871 & 0.889 \\
    Hounslow               & 651  & 1.054 & 0.808 & 0.822 & 0.832 \\
    Islington              & 349  & 0.787 & 0.797 & 0.792 & 0.858 \\
    Kensington and Chelsea & 260  & 0.747 & 0.759 & 0.760 & 0.811 \\
    Kingston upon Thames   & 391  & 1.209 & 0.835 & 0.838 & 0.875 \\
    \textbf{Lambeth}       & 577  & 0.862 & 0.828 & 0.823 & \textbf{0.822} \\
    Lewisham               & 609  & 0.986 & 0.808 & 0.809 & 0.828 \\
    \textbf{Merton}        & 457  & 1.035 & 0.797 & 0.788 & \textbf{0.793} \\
    Newham                 & 465  & 0.832 & 0.796 & 0.785 & 0.820 \\
    Redbridge              & 539  & 1.102 & 0.795 & 0.814 & 0.838 \\
    Richmond upon Thames   & 484  & 1.317 & 0.800 & 0.808 & 0.838 \\
    Southwark              & 583  & 0.860 & 0.817 & 0.824 & 0.845 \\
    Sutton                 & 421  & 1.487 & 0.828 & 0.824 & 0.849 \\
    Tower Hamlets          & 381  & 0.949 & 0.812 & 0.882 & 0.828 \\
    Waltham Forest         & 492  & 0.888 & 0.811 & 0.793 & 0.856 \\
    Wandsworth             & 481  & 0.815 & 0.792 & 0.803 & 0.874 \\
    Westminster            & 419  & 0.763 & 0.777 & 0.773 & 0.781 \\
    \bottomrule
  \end{tabular}
\end{table}

The boroughs where GATv2 wins---Hackney, Barnet, City of London, Haringey,
Lambeth, Merton, and Bromley---are concentrated in inner and central London.
These are areas characterised by dense, overlapping bus route networks where
multiple routes serve the same corridors. In such environments, route-connectivity
edges carry genuine signal: a stop on a busy through-route benefits from
the demand patterns of the preceding and following stops on the same service.
RF, operating on tabular features without any relational structure, cannot
exploit this information.

By contrast, in outer London boroughs such as Havering, Hillingdon, Barking,
and Ealing, RF outperforms GATv2 by wider margins. These areas have sparser
and less overlapping route structures, and the demand variation between stops
is explained more by land-use accessibility differences than by network
topology. In such contexts, the graph adds noise rather than signal.

\paragraph{City of London: a resolved anomaly.}
An earlier version of the pipeline, which applied a global StandardScaler
fitted on the full dataset before splitting (a data leakage error), produced
a WMAPE of 9.36 for City of London. This was traced to the extreme employment
accessibility values of the City (the financial district has employment
accessibility 10--40 times the outer-borough median), which caused the scaler
to produce distorted inputs at evaluation time. After correcting the scaler to
be fitted per fold on training data only, and applying output clamping to
$[0,\, \log(1 + 2 \cdot y_{\text{max}}^{\text{train}})]$, the City of London
WMAPE reduced to 0.798, competitive with RF (0.819). This episode illustrates
the importance of correct cross-validation implementation in spatially
heterogeneous datasets.


% ============================================================
\chapter{Discussion}
\label{ch:discussion}

\section{What the graph adds and when}
The central finding of this dissertation is that graph structure---in the
specific form of route-connectivity edges combined with geographic KNN
edges---provides a measurable but context-dependent benefit for cold-start
bus demand prediction. GATv2 outperforms RF in boroughs where the bus
network is sufficiently dense for route-level passenger flow to be a
discriminative signal. In sparser outer-London networks, the same edges
become less informative, and the tabular feature model performs better.

This result is consistent with the theoretical properties of message passing:
the operation is only beneficial when the aggregated neighbourhood information
differs meaningfully from the node's own representation
\citep{hamilton2017inductive}. In this dataset, 96.4\% of stops share their
LSOA, and therefore their AI23 feature vector, with at least one neighbour.
The addition of geographic coordinates partially restores node uniqueness, but
the residual homophily remains the primary constraint on GATv2's advantage.

\section{Comparison with Zheng et al. (2025)}
\citet{zheng2025mfstgat} report GATv2 outperforming Random Forest in their
setting, which contrasts with the aggregate result found here. The most
likely explanation is the difference in feature granularity: Zheng et al.
use stop-level POI features (counts of nearby restaurants, offices, shops,
and other amenities within a fixed buffer radius) that are unique to each
individual stop. This ensures that message passing aggregates genuinely
diverse representations from heterogeneous neighbours. The LSOA-level
features used here cannot replicate this diversity within the same
administrative zone.

A secondary factor is evaluation design. It is not clear from \citet{zheng2025mfstgat}
whether their evaluation uses spatially blocked folds or random splits.
Random splits, by allowing nearby stops to appear in both training and test
sets, tend to inflate estimates of GNN performance relative to tabular
baselines, because the GNN can effectively interpolate over its observed
neighbours.

\section{Practical implications for TfL}
Despite trailing RF in aggregate, GATv2 offers a practical argument for
deployment in specific contexts. For a new stop being added to an existing
high-frequency inner-London corridor---for example, an extension of a route
into a new development in Hackney or Barnet---GATv2 can leverage the demand
patterns of the existing stops on the same route, yielding predictions that
RF, constrained to accessibility features alone, cannot match.

For new stops in areas with limited connectivity, such as new settlements
at the periphery of the network, RF is the safer choice, being more stable
across the full distribution of boroughs (standard deviation 0.026 vs 0.032
for GATv2). A practical deployment strategy might therefore select between
the two models based on the route density of the target area.

\section{Limitations}
\label{sec:limitations}
Several limitations should be noted. The BUSTO demand target is a
survey-derived estimate of a typical autumn weekday rather than a verified
count of individual passengers; the fractional boarding values reflect
statistical weighting. The AI23 accessibility features are at LSOA level,
and all stops within the same zone receive identical feature vectors, which
limits the information available to message passing. The cold-start scenario
is simulated through geographic hold-out rather than observed on an actual
new development, so the work demonstrates geographic transferability rather
than true deployment performance. The model predicts a typical-day static
average and does not capture the temporal dynamics (AM peak, PM peak, evening)
that a more complete model would require. Results are specific to London; the
degree to which the findings generalise to other cities with different network
densities and land-use structures is not established by this study.


% ============================================================
\chapter{Conclusions}
\label{ch:conclusion}

\section{Summary of findings}
This dissertation investigated whether an inductive GATv2 can predict bus-stop
passenger demand in geographically unseen areas---the cold-start problem---using
accessibility indicators and route-connectivity as the sole inputs. The model
was evaluated under 33-fold leave-borough-out cross-validation across all London
boroughs, producing the following headline results:

\begin{itemize}
  \item GATv2 (WMAPE 0.838) and Random Forest (WMAPE 0.812) both substantially
  outperform the Historical Average baseline (WMAPE 1.082), confirming that
  accessibility features carry meaningful signal for demand prediction in the
  absence of historical ridership.
  \item GATv2 outperforms RF in 7 of 33 boroughs, concentrated in inner London
  where dense route networks make route-connectivity edges informative.
  \item The MLP ablation (WMAPE 0.814) is essentially equivalent to RF,
  confirming that node features account for the majority of predictive power
  and that graph structure contributes a bounded, geography-dependent increment.
  \item The 0.026 WMAPE gap between RF and GATv2 is attributed to LSOA-level
  feature homophily: 96.4\% of stops share identical AI23 feature vectors,
  limiting the diversity of information that message passing can aggregate.
\end{itemize}

\section{Limitations}
The principal limitations of this work are the LSOA-level granularity of
accessibility features, the simulated rather than observed cold-start
scenario, the absence of temporal dynamics in the demand target, and the
restriction to London. Each of these points directly towards future work.

\section{Future work}
The most impactful immediate extension is the addition of stop-level POI
features derived from OpenStreetMap via the Overpass API. Querying counts of
nearby restaurants, offices, retail outlets, and other amenities within a
fixed buffer (e.g., 400 metres) around each stop would give every node a
unique feature vector regardless of LSOA membership, addressing the primary
constraint on GATv2's relative performance identified in this study.

A second direction is the incorporation of temporal dynamics. The BUSTO dataset
contains 15-minute interval data; restructuring the target as a time-series
of boardings across the operating day would enable spatio-temporal GNN
architectures such as MF-STGAT \citep{zheng2025mfstgat} to be evaluated in
the inductive cold-start setting defined here. This would make the comparison
with the state of the art in transit demand modelling more direct.

A third direction concerns prescriptive rather than descriptive modelling:
using the demand predictions from this work as inputs to a route optimisation
algorithm (e.g., multi-objective evolutionary methods such as NSGA-II) to
suggest where new stops should be placed in a developing area, given predicted
demand and network cost constraints.

Finally, generalising the pipeline to cities beyond London would establish
whether the finding that GATv2's advantage is concentrated in high-density
inner-city networks is a general property of bus demand prediction or an
artefact of London's specific topology.


% ============================================================
\bibliographystyle{plainnat}
\bibliography{references}

\end{document}
